Mỗi CPU có một kernel thread riêng chạy \texttt{scheduler()} (\texttt{kernel/proc.c}) trong một vòng lặp vô hạn:
\begin{enumerate}
    \item Bật ngắt (\texttt{intr\_on()}), để tránh deadlock khi mọi process đều đang chờ ngắt thiết bị.
    \item Duyệt \texttt{proc[]} theo vòng tròn (round-robin). Với mỗi \texttt{p}, lấy \texttt{p->lock}; nếu \texttt{p} đang \texttt{RUNNABLE} thì đặt \texttt{RUNNING}, ghi \texttt{c->proc = p}, tăng \texttt{p->nsched} (phần nhóm thêm), rồi \texttt{swtch(\&c->context, \&p->context)} để chạy \texttt{p}.
    \item Khi \texttt{p} nhường CPU, \texttt{swtch} quay về đây: đặt \texttt{c->proc = 0}, nhả \texttt{p->lock}, xét process tiếp theo.
    \item Duyệt hết bảng mà không có process nào để chạy thì thực hiện lệnh \texttt{wfi} để chờ ngắt.
\end{enumerate}

\textbf{\texttt{sched()}} là đường duy nhất để một process trả CPU về scheduler: nó gọi \texttt{swtch(\&p->context, \&mycpu()->context)}, có lưu và khôi phục \texttt{intena} của CPU. Trước đó hàm kiểm tra bốn điều kiện, sai điều nào thì \texttt{panic}:
\begin{itemize}
    \item Phải đang giữ \texttt{p->lock}, để CPU khác không chọn chạy \texttt{p} khi \texttt{p} chưa lưu xong context.
    \item Chỉ giữ đúng một khóa (\texttt{noff == 1}). Nhường CPU khi còn giữ khóa khác sẽ làm các process cần khóa đó phải quay chờ mãi.
    \item \texttt{state} đã khác \texttt{RUNNING}: người gọi phải đặt trạng thái mới trước (\texttt{RUNNABLE}, \texttt{SLEEPING} hoặc \texttt{ZOMBIE}).
    \item Ngắt đang tắt, để không có ngắt nào chen vào giữa lúc chuyển context.
\end{itemize}

\textbf{\texttt{yield()}} lấy \texttt{p->lock}, đặt \texttt{RUNNABLE} rồi gọi \texttt{sched()}; khi được chạy lại thì nhả khóa. \texttt{usertrap()} và \texttt{kerneltrap()} gọi \texttt{yield()} mỗi khi gặp ngắt timer (\texttt{which\_dev == 2}), nên một process chạy vòng lặp vô hạn vẫn bị lấy lại CPU sau mỗi tick.

\textbf{\texttt{p->lock} được giữ qua \texttt{swtch}.} Khóa này được lấy ở một thread và nhả ở thread khác. Scheduler lấy \texttt{p->lock} rồi chuyển sang \texttt{p}, và chính \texttt{p} nhả khóa (trong \texttt{yield()}, \texttt{sleep()}, hoặc \texttt{forkret()} nếu là lần chạy đầu). Ngược lại, \texttt{p} lấy khóa trước khi gọi \texttt{sched()}, và scheduler nhả nó sau khi \texttt{swtch} quay về. Nhờ vậy trạng thái của \texttt{p} được bảo vệ suốt quá trình chuyển, không CPU nào chen vào giữa được.
